
%% ================================================================
%%  5.  EXPERIMENTAL DESIGN
%% ================================================================
\section{Experimental Design}\label{sec:design}

\subsection{Evaluation Protocol}

We define a rigorous evaluation protocol to ensure fair, reproducible
comparison across all agents:

\begin{itemize}
  \item Each of the \textbf{32 scenarios} is evaluated over
        \textbf{30 independent episodes} with distinct random seeds
        (seeds 1000--1029), all of which are unseen during RL training.
  \item All agents---classical OR methods, heuristics, and RL
        agents---are evaluated on \textbf{identical demand
        realizations} within each seed, ensuring that performance
        differences reflect algorithmic capability rather than
        stochastic variation.
  \item RL agents are evaluated in \textbf{deterministic mode} (mean
        of the Gaussian policy, no exploration noise).
\end{itemize}

\subsection{Performance Metrics}

We report the following metrics for each agent--scenario pair:

\paragraph{Total Episode Profit ($\Pi$).}
The primary metric. Computed as $\Pi = \sum_{t=0}^{T-1} r_t$ where
$r_t$ is the per-period net profit (\Cref{eq:reward}) and $T = 30$.
We report the mean $\bar{\Pi}$ and standard deviation $\sigma_\Pi$
over 30~episodes.

\paragraph{Fill Rate.}
The fraction of demand fulfilled:
\begin{equation}\label{eq:fill_rate}
  \text{FR} = \frac{\sum_{t} S_t}{\sum_{t} D_t},
\end{equation}
where $S_t$ is units sold and $D_t$ is realized demand. Fill rate
captures customer service quality independent of the cost structure.

\paragraph{Average On-Hand Inventory.}
$\bar{X} = \frac{1}{T} \sum_{t} \sum_{n} X_t^n$, measuring capital
tied up in inventory.

\paragraph{Total Unfulfilled Demand.}
$U_{\text{total}} = \sum_{t} U_t$, measuring cumulative stock-outs.

\paragraph{Wall-Clock Computation Time.}
Average time per scenario (30 episodes), measured on a MacBook~2019
(Intel Core~i5, 8~GB~RAM). This enables practitioners to assess the
speed--quality trade-off for real-time applications.

\subsection{Statistical Testing}

To distinguish genuine performance differences from stochastic noise,
we employ \textbf{paired $t$-tests} between agent pairs. Since each
agent is evaluated on the same 30~seeds per scenario, the paired test
controls for demand variability. We report $p$-values and flag
differences as significant at $\alpha = 0.05$.

\subsection{Information-Value Metrics}

We compute two classical stochastic programming metrics that quantify
the economic value of information:

\paragraph{Value of Perfect Information (VPI).}
\begin{equation}\label{eq:vpi}
  \text{VPI} = \Pi^{\text{Oracle}} - \Pi^{\text{MSSP (Informed)}},
\end{equation}
the profit premium from knowing future demand perfectly. VPI
represents the maximum amount a decision-maker should pay for perfect
demand forecasts.

\paragraph{Value of the Stochastic Solution (VSS).}
\begin{equation}\label{eq:vss}
  \text{VSS} = \Pi^{\text{MSSP}} - \Pi^{\text{DLP}},
\end{equation}
the profit premium from modelling demand uncertainty (stochastic
program) versus using a deterministic approximation. A high VSS
indicates that demand uncertainty is economically significant---the
extra computational cost of MSSP is justified.

\paragraph{Optimality Gap.}
The relative optimality gap for each agent:
\begin{equation}\label{eq:gap}
  \text{Gap}(a) = \frac{\Pi^{\text{Oracle}} - \Pi^{a}}
                       {\Pi^{\text{Oracle}}} \times 100\%.
\end{equation}
This normalizes performance relative to the theoretical upper bound,
enabling comparison across scenarios with different absolute profit
levels.

\subsection{Reproducibility}

To support reproducibility, all code, trained models, configuration
files, and evaluation scripts are released in the companion
repository. The evaluation pipeline is fully automated: a single
command runs all 10~agents across all 32~scenarios with the prescribed
seeds and produces CSV results files, summary statistics, and
publication-ready figures.
